% =====================================================================
%  CAPÍTULO 7 — RESULTADOS Y DISCUSIÓN
% =====================================================================

\chapter{Resultados y discusión}
\label{cap:resultados}

\section{Organización de la presentación}
\label{sec:organizacion-resultados}

Los resultados se presentan en el orden en que se obtuvieron, y no en el orden lógico del plan original. Ese orden es en sí mismo un resultado: el hallazgo del primer bloque invalidó la hipótesis de partida y obligó a reordenar los bloques posteriores. Presentarlos en el orden previsto daría la imagen falsa de un diseño cerrado desde el principio.

Salvo indicación expresa, las cifras corresponden a la tasa de error por carácter con decodificación CTC voraz, calculada de forma micro sobre la partición de validación completa (\cref{sec:protocolo-evaluacion}). Cuando se comparan objetivos de salida distintos, el eje común es la tasa de error por palabra, conforme a la \cref{sec:metricas-teoria}.

Antes de los bloques conviene registrar la comprobación de capacidad. Entrenado sobre un subconjunto de 348 ensayos hasta el sobreajuste deliberado, el sistema alcanza una tasa de error por carácter de $0{,}052$. La cadena completa es, por tanto, capaz de aprender la tarea: cuanto se observe después es un problema de generalización y no de conexión entre la señal y el modelo.

\section{Bloques A y H: la aportación del preentrenamiento}
\label{sec:resultados-A}

\subsection{Resultados}

\begin{table}[H]
  \centering
  \caption{Aportación del preentrenamiento. Tasa de error por carácter mínima sobre validación, época del mínimo, presupuesto y coste en unidades de cómputo.}
  \label{tab:bloque-a}
  \small
  \begin{tabular}{llcccc}
    \toprule
    \textbf{Ejec.} & \textbf{Codificador} & \textbf{CER} & \textbf{Época} & \textbf{Épocas} & \textbf{Coste} \\
    \midrule
    \multicolumn{6}{l}{\textit{Inicialización aleatoria, ajuste completo}} \\
    H1 & Semilla 2 & \textbf{0,297} & 40 & 40 & 8,8 \\
    H2 & Semilla 3 & 0,314 & 40 & 40 & 8,8 \\
    A1 & Semilla 1 & 0,324 & 28 & 30 & \textbf{6,6} \\
    \midrule
    \multicolumn{6}{l}{\textit{Pesos de OWSM}} \\
    v5 & Ajuste completo & 0,362 & 60 & 60 & 15,2 \\
    v3 & Ajuste completo & 0,381 & 37 & 38 & 8,9 \\
    D0 & Congelado, presupuesto ampliado & 0,440 & 70 & 70 & 14,6 \\
    A2 & Congelado & 0,658 & 29 & 30 & 6,3 \\
    H3 & Ajuste completo, tasa cuadruplicada & 0,689 & 39 & 40 & 8,7 \\
    \midrule
    \multicolumn{6}{l}{\textit{Controles adicionales}} \\
    A4 & Aleatorio, congelado & 0,414 & 34 & 40 & 8,8 \\
    A1$_{\text{lr}}$ & Aleatorio, tasa unificada & 0,854\rlap{$^{*}$} & --- & 30 & 6,5 \\
    \bottomrule
  \end{tabular}

  \vspace{0.4em}
  {\footnotesize $^{*}$ Ejecución divergente. Véase la \cref{sec:divergencia}.}
\end{table}

Las tres semillas con inicialización aleatoria promedian $0{,}312$ con una desviación típica de $0{,}014$ y un rango de $[0{,}297,\ 0{,}324]$. La mejor ejecución con pesos de OWSM y receta comparable alcanza $0{,}381$, situada a \textbf{5,1 desviaciones típicas} por encima de la media aleatoria. La diferencia excede con holgura tanto la dispersión entre semillas como el intervalo de confianza de la medida establecido en la \cref{sec:precision-medida}.

\begin{figure}[H]
  \centering
  \fbox{\parbox{0.9\textwidth}{\vspace{2.3cm}\centering\textit{[FIGURA: curvas de CER en validación frente a época para A1, H1 y H2 (aleatorias, en un color), v3 y v5 (OWSM, en otro), y A2/D0 (congelado). Es la figura principal de la memoria: debe verse la banda de las tres semillas aleatorias enteramente por debajo de las preentrenadas.]}\vspace{2.3cm}}}
  \caption{Evolución de la tasa de error por carácter según la inicialización del codificador.}
  \label{fig:bloque-a-curvas}
\end{figure}

\subsection{Las dos objeciones previsibles, resueltas}

El bloque H se diseñó para anticipar las dos objeciones evidentes a un resultado que contradice la literatura previa. Ambas quedan resueltas en contra de la objeción.

\textbf{¿Es reproducible?} Sí. Las tres semillas ocupan un intervalo de 2,7 puntos porcentuales, y la peor de ellas sigue siendo 5,7 puntos mejor que la mejor ejecución preentrenada de receta comparable. El efecto no es la fluctuación de una única ejecución afortunada.

\textbf{¿Se le dio a OWSM su mejor oportunidad?} Sí, y empeora. La hipótesis natural era que la tasa de aprendizaje reducida del codificador impidiese a los pesos preentrenados alejarse lo suficiente de su punto de partida, penalizando por construcción a la condición preentrenada. Al cuadruplicar esa tasa el resultado no mejora, sino que \textbf{se degrada de $0{,}381$ a $0{,}689$}. Lejos de estar infra-adaptados, los pesos de OWSM se destruyen en cuanto se les permite moverse con libertad, sin que el modelo recupere por ello el rendimiento de una inicialización aleatoria. La objeción queda refutada empíricamente.

\subsection{El control que cierra el argumento}
\label{sec:control-a4}

El resultado más contundente del bloque no es ninguna de las dos comparaciones anteriores, sino el control adicional A4: un codificador \textbf{inicializado aleatoriamente y congelado}, esto es, una proyección aleatoria fija sobre la que sólo se entrenan el adaptador y la cabeza CTC.

Ese control alcanza $0{,}414$ en 40 épocas, frente a los $0{,}440$ que obtiene el codificador de OWSM congelado tras 70 épocas, casi el doble de presupuesto. Dicho de otro modo: \textbf{como extractor de características fijo, un codificador aleatorio sin entrenar resulta al menos tan útil como uno preentrenado sobre 180.000 horas de habla}, y lo consigue con menos cómputo.

Este control transforma una conclusión débil en una fuerte. No se trata únicamente de que el preentrenamiento no aporte información aprovechable: es que las representaciones que induce no superan a una proyección aleatoria, lo que sitúa su contenido transferible a esta modalidad en un valor indistinguible de cero.

\subsection{La tasa de aprendizaje diferenciada es un requisito de estabilidad}
\label{sec:divergencia}

Tres ejecuciones independientes con tasa unificada entre codificador y componentes nuevos divergieron: la de inicialización aleatoria ($0{,}854$) y las dos de descongelado parcial ($0{,}848$ y $0{,}859$). En todas ellas la pérdida de entrenamiento se estanca mientras la de validación asciende, patrón característico de divergencia y no de entrenamiento lento.

Que el fenómeno se reproduzca tanto con codificador aleatorio como con codificador preentrenado confirma lo anticipado en la \cref{sec:estrategias-adaptacion}: la tasa diferenciada no cumple aquí la función de proteger unos pesos preentrenados, sino la de estabilizar el acoplamiento entre un adaptador que aprende deprisa y un codificador profundo cuyas representaciones aún no están organizadas.

\subsection{Contraste con la literatura previa}
\label{sec:contraste-fiedler}

El resultado contradice a Fiedler \textit{et al.} \cite{fiedler2025wav2vec2brain}, que aplicando el mismo control sobre wav2vec~2.0 y los datos del participante T12 obtuvieron una ventaja de más de veinte puntos porcentuales \textit{a favor} de los pesos preentrenados.

Con la evidencia ahora disponible, dos de las cuatro explicaciones que cabía plantear quedan descartadas. La restricción impuesta por la tasa de aprendizaje se descarta por el resultado de H3. Y la hipótesis de que un adaptador demasiado potente estuviese absorbiendo la función del codificador preentrenado se descarta por el bloque C (\cref{sec:resultados-C}): al reducir la capacidad del adaptador el sistema empeora de forma monótona, sin que la ventaja del preentrenamiento reaparezca en ningún punto.

Quedan en pie dos explicaciones, ambas relativas a \textit{qué} se preentrena y no a \textit{cómo} se adapta:

\begin{enumerate}[label=\textbf{H\arabic*.}, leftmargin=*, itemsep=0.3em]
  \item \textbf{Naturaleza del preentrenamiento.} wav2vec~2.0 se preentrena de forma autosupervisada sobre forma de onda cruda mediante un extractor convolucional, mientras que OWSM lo hace con supervisión débil sobre banco de filtros mel. Un codificador construido sobre una representación espectral estaría especializado en una estructura de la que la señal neuronal carece por completo, mientras que las representaciones autosupervisadas de bajo nivel serían más agnósticas a la modalidad. La hipótesis predice que la propiedad determinante no es haber visto habla, sino el nivel de abstracción al que se aprendió a representarla.
  \item \textbf{Diferencias del conjunto de datos.} El horizonte de 20 meses del corpus empleado aquí, frente a los cuatro meses del utilizado por ellos, impone una carga de no estacionariedad muy superior, que podría dominar sobre cualquier ventaja de inicialización.
\end{enumerate}

Ninguna de las dos es contrastable con los recursos de este trabajo, y ambas se recogen como línea de continuación en el \cref{cap:limitaciones}.

\section{Bloque D: la curva de conservación}
\label{sec:resultados-D}

\begin{table}[H]
  \centering
  \caption{Rendimiento en función de la cantidad de codificador preentrenado que se adapta. La pendiente final es la del descenso de la tasa de error en la última época del presupuesto.}
  \label{tab:bloque-d}
  \small
  \begin{tabular}{llccc}
    \toprule
    \textbf{Ejec.} & \textbf{Condición} & \textbf{CER} & \textbf{Épocas} & \textbf{Pendiente final} \\
    \midrule
    A2 & OWSM congelado                & 0,658 & 30 & $-0{,}022$ \\
    D0 & OWSM congelado, ampliado      & 0,440 & 70 & $-0{,}001$ \\
    D2 & OWSM, 2 capas superiores      & 0,521 & 30 & $-0{,}008$ \\
    D3 & OWSM, 4 capas superiores      & 0,493 & 30 & $-0{,}008$ \\
    D1 & OWSM congelado con LoRA       & 0,403 & 30 & $-0{,}004$ \\
    v3 & OWSM, ajuste completo         & 0,381 & 38 & $-0{,}004$ \\
    A1 & Aleatorio, ajuste completo    & \textbf{0,324} & 30 & $-0{,}001$ \\
    \bottomrule
  \end{tabular}
\end{table}

La curva \textbf{no es monótona}, y su interpretación exige atender a la columna de pendientes. Las dos condiciones de descongelado parcial conservan al final del presupuesto una pendiente de $-0{,}008$ por época, un orden de magnitud superior a la de las condiciones convergidas: no han alcanzado su meseta y sus cifras no representan su rendimiento asintótico. El mismo argumento explica el salto del codificador congelado, de $0{,}658$ a $0{,}440$, al más que duplicar su presupuesto.

Descontando las condiciones no convergidas, el orden que sí puede afirmarse es que \textbf{el ajuste completo supera a la adaptación de rango bajo, esta supera a la congelación, y la inicialización aleatoria supera a todas}. Resulta destacable que la adaptación de rango bajo sobre codificador congelado ($0{,}403$) mejore en 3,7 puntos a la congelación simple pese a disponer de menos de la mitad de presupuesto ($0{,}440$ con 70 épocas), lo que confirma que su beneficio procede de la restricción del subespacio de adaptación y no de la capacidad añadida, conforme al argumento de la \cref{sec:lora-teoria}.

Lo que la curva no muestra en ningún punto es un óptimo intermedio: no existe una cantidad de preentrenamiento conservado que resulte preferible a prescindir de él por completo.

\section{Bloque F: el tamaño de la arquitectura}
\label{sec:resultados-F}

Establecido que los pesos no aportan, queda por determinar si lo hace el tamaño de la arquitectura que los albergaba.

\begin{table}[H]
  \centering
  \caption{Efecto de la profundidad del codificador con inicialización aleatoria.}
  \label{tab:bloque-f}
  \small
  \begin{tabular}{lcccc}
    \toprule
    \textbf{Bloques} & \textbf{CER} & \textbf{WER voraz} & \textbf{Épocas} & \textbf{Minutos} \\
    \midrule
    6 (arquitectura de OWSM) & 0,324 & 0,571 & 30 & 232 \\
    3                        & 0,334 & 0,569 & 40 & 272 \\
    2                        & 0,353 & 0,602 & 40 & 259 \\
    \bottomrule
  \end{tabular}
\end{table}

Reducir el codificador a la mitad cuesta un punto porcentual de tasa de error por carácter y deja la tasa de error por palabra prácticamente inalterada. Reducirlo a un tercio cuesta tres puntos. La degradación es, por tanto, muy suave, lo que indica que \textbf{tampoco el tamaño de la arquitectura de OWSM está justificado por la tarea}: el problema admite codificadores sensiblemente menores sin pérdida apreciable.

Una precisión honesta sobre esta tabla: las ejecuciones reducidas se entrenaron 40 épocas frente a las 30 de la completa, de modo que el ahorro no se materializa en la columna de tiempo total. El ahorro real está en el coste por época, y una comparación limpia de eficiencia exigiría igualar también el presupuesto.

\section{Bloque B: granularidad de la unidad de salida}
\label{sec:resultados-B}

\subsection{Resultados}

\begin{table}[H]
  \centering
  \caption{Comparación de unidades de salida con decodificación voraz. La columna de error por palabra es el único eje comparable entre filas.}
  \label{tab:bloque-b}
  \small
  \begin{tabular}{lrccc}
    \toprule
    \textbf{Unidad} & \textbf{Clases} & \textbf{CER} & \textbf{WER} & \textbf{Épocas} \\
    \midrule
    Palabra              & Cerrado & \textbf{0,222} & \textbf{0,282} & 40 \\
    Subpalabra de OWSM   & 50.002  & 0,247 & 0,323 & 20 \\
    Subpalabra propia    & 256     & 0,251 & 0,424 & 40 \\
    Carácter             & 29      & 0,267 & 0,535 & 40 \\
    Fonema               & ---     & \multicolumn{2}{c}{No produce texto} & 40 \\
    \bottomrule
  \end{tabular}
\end{table}

El resultado es monótono y \textbf{contrario a la predicción formulada en los \cref{cap:datos,cap:arquitectura}}: cuanto más gruesa es la unidad de salida, mejor es la transcripción final. La diferencia entre el mejor y el peor objetivo que produce texto asciende a 25 puntos porcentuales de tasa de error por palabra, muy por encima de cualquier otro efecto medido en el trabajo, incluida la contribución del preentrenamiento.

En particular, el vocabulario de subpalabras de OWSM, que esta memoria había descartado \textit{a priori} por escasez de ejemplos por clase, resulta ser el segundo mejor objetivo, y lo consigue con la mitad de presupuesto que los demás y sin haber alcanzado su meseta: su mínimo se produce en la penúltima época con una pendiente todavía de $-0{,}007$ por época.

\subsection{Por qué falló la predicción}

El razonamiento original era que repartir poco más de ocho mil frases entre 50.002 clases dejaría a la inmensa mayoría sin ejemplos, y que un vocabulario reducido sería por tanto preferible. El razonamiento es correcto en los hechos y erróneo en la consecuencia, por dos motivos que la literatura de reconocimiento del habla tenía documentados y que no se tuvieron en cuenta.

El primero es que \textbf{en CTC las clases no observadas no penalizan}. Una clase que nunca aparece en el entrenamiento simplemente no se predice: no compite por masa de probabilidad ni degrada a las demás. El coste de un vocabulario grande es de memoria y de tiempo de cálculo de la función \textit{softmax}, no de calidad.

El segundo, y determinante, es el efecto de la unidad sobre la \textbf{longitud de la secuencia objetivo}. Frente a los caracteres, las subpalabras producen secuencias de salida más cortas y reducen así la dificultad de modelar la dependencia entre salidas, algo especialmente importante en modelos basados en CTC por su hipótesis de independencia condicional \cite{higuchi2021hierarchical}. El fenómeno es conocido: con unidades grandes como la palabra, las limitaciones de CTC se sortean en parte, mientras que con unidades pequeñas como el carácter se vuelven prominentes \cite{das2018advancing}. En reconocimiento del habla convencional, sustituir el objetivo de carácter por subpalabras produce reducciones relativas de tasa de error del orden del \SI{13}{\percent} sin diccionario ni modelo de lenguaje externo \cite{zeyer2018subword}.

A ello se añade una circunstancia específica de este montaje. Con ventanas de \SI{20}{\milli\second} y submuestreo por dos, la secuencia de entrada queda a una trama cada \SI{40}{\milli\second}. CTC exige que la entrada sea más larga que la etiqueta, y con objetivo de carácter esa holgura se estrecha considerablemente. Es un problema documentado al reducir la resolución temporal, y su solución habitual consiste precisamente en adoptar subpalabras en lugar de caracteres \cite{chen2023fasthubert}.

La conclusión metodológica es que \textbf{el argumento de los ejemplos por clase, aplicado sin considerar la longitud de la secuencia objetivo, conduce a la decisión equivocada}. Es un error de razonamiento a priori que sólo el experimento podía corregir, y constituye por sí mismo un argumento a favor de contrastar empíricamente las decisiones de diseño en lugar de deducirlas.

\subsection{Matices y limitaciones del hallazgo}

Tres precisiones impiden leer la tabla como una recomendación simple de emplear siempre la unidad más gruesa.

El objetivo de \textbf{palabra opera sobre vocabulario cerrado} y es, por construcción, incapaz de emitir una palabra no vista en entrenamiento. Dado que la partición de validación incorpora deliberadamente corpus ausentes del entrenamiento (\cref{sec:b2t25}), su ventaja está inflada respecto a lo que ofrecería en uso real, donde el vocabulario abierto es un requisito. Se trata de una limitación estructural y no de una cuestión de volumen de datos: la literatura señala precisamente la ausencia del problema de palabras fuera de vocabulario como la ventaja de caracteres y subpalabras frente a palabras \cite{zeyer2018subword}.

El objetivo de \textbf{fonema no produce texto ortográfico}, de modo que no puede situarse en el eje común sin un decodificador con léxico. Su ausencia de la comparación es una limitación reconocida y resulta especialmente lamentable, ya que es la unidad que emplea la práctica totalidad de la literatura del campo y la que el argumento neurocientífico señalaba como natural.

Por último, la comparación se realiza con \textbf{decodificación voraz y sin modelo de lenguaje externo}. Es previsible que la fusión con un modelo de lenguaje reduzca la ventaja de las unidades gruesas, puesto que buena parte de lo que aportan es precisamente la estructura lingüística implícita que el modelo de lenguaje suministra por otra vía.

\subsection{Respuesta a PI2}

La corteza motora del habla puede leerse a granularidades muy distintas, pero \textbf{la unidad que minimiza el error del texto final no es la más próxima a la articulación sino la más próxima al lenguaje}. El factor determinante no es la plausibilidad neurocientífica de la unidad, sino la longitud de la secuencia que induce y la cantidad de estructura lingüística que incorpora implícitamente, compensando la hipótesis de independencia condicional de CTC.

\section{Bloque E: generalización a sesiones no vistas}
\label{sec:resultados-E}

\begin{table}[H]
  \centering
  \caption{Rendimiento sobre las cinco sesiones excluidas del entrenamiento, según el tratamiento de la capa de sesión.}
  \label{tab:bloque-e}
  \small
  \begin{tabular}{lccc}
    \toprule
    \textbf{Tratamiento} & \textbf{CER} & \textbf{WER voraz} & \textbf{Época del mínimo} \\
    \midrule
    Sin capa de sesión               & \textbf{0,593} & 0,848 & 33 \\
    Matriz en la identidad           & 0,747 & 0,984 & 10 \\
    Media de las matrices entrenadas & 0,752 & 0,970 & 12 \\
    \bottomrule
  \end{tabular}
\end{table}

El resultado es inequívoco: ante un día de registro no visto, \textbf{el mejor tratamiento de la capa de sesión es no tenerla}. Prescindir de ella mejora en más de 15 puntos porcentuales sobre cualquiera de las dos estrategias de sustitución, y ni la identidad ni el promedio de las matrices entrenadas ofrecen ventaja apreciable la una sobre la otra.

La columna de la época del mínimo revela el mecanismo. Las dos configuraciones con capa de sesión alcanzan su mejor rendimiento sobre días no vistos en las épocas 10 y 12, y a partir de ahí se degradan pese a que su rendimiento sobre días vistos continúa mejorando. La capa de sesión no está aprendiendo una corrección generalizable de la deriva del implante, sino \textbf{memorizando cada día concreto}, y esa memorización resulta progresivamente perjudicial fuera de los días memorizados.

Contrastado con el resto del trabajo, el compromiso es severo. Con todos los días disponibles y capa de sesión, el sistema alcanza $0{,}324$; sobre días no vistos, la mejor configuración se sitúa en $0{,}593$. La degradación de casi 27 puntos porcentuales es, con diferencia, la mayor medida en todo el trabajo, y sitúa la no estacionariedad, y no la elección de codificador, como el problema dominante del despliegue real. Es también la conclusión de mayor relevancia clínica, por corresponder al escenario del primer día tras el implante o de la primera sesión tras un periodo sin uso.

Los aproximadamente doce millones de parámetros que cuesta la capa de sesión compran, en definitiva, rendimiento dentro de la distribución a cambio de capacidad de generalización fuera de ella.

\section{Bloque C: capacidad del módulo adaptador}
\label{sec:resultados-C}

\begin{table}[H]
  \centering
  \caption{Comparación de variantes de adaptador. Las tres comparten factor de submuestreo dos.}
  \label{tab:bloque-c}
  \small
  \begin{tabular}{lccc}
    \toprule
    \textbf{Variante} & \textbf{CER} & \textbf{WER voraz} & \textbf{Minutos} \\
    \midrule
    Profunda            & \textbf{0,324} & 0,571 & 232 \\
    Convolucional 1D    & 0,399 & 0,643 & 244 \\
    Lineal              & 0,475 & 0,710 & 185 \\
    \bottomrule
  \end{tabular}
\end{table}

La advertencia formulada en la \cref{sec:diseno-experimental} sobre una posible confusión entre arquitectura del adaptador y resolución temporal \textbf{no llegó a materializarse}: las tres variantes se ejecutaron con factor de submuestreo dos, de modo que la comparación es limpia.

El adaptador profundo justifica su coste con claridad, mejorando en 7,5 puntos porcentuales a la variante convolucional y en 15,1 a la lineal, con un tiempo por época comparable al de la primera. La separación explícita entre mezcla espacial de electrodos y convolución temporal, argumentada en la \cref{sec:frontend-deep}, se traduce en una diferencia sustancial de rendimiento.

Este bloque cumple además una función no prevista en el plan original, la de contrastar la hipótesis de que un adaptador excesivamente potente estuviese absorbiendo la función del codificador preentrenado. La predicción de esa hipótesis era que la ventaja del preentrenamiento reapareciese al reducir la capacidad del adaptador. No ocurre: al debilitarlo, el sistema empeora de forma monótona y en ningún punto la condición preentrenada resulta preferible.

\section{Decodificación con modelo de lenguaje}
\label{sec:resultados-lm}

% TODO (Kiril): esta seccion sigue con cifras preliminares. En los CSV las
% columnas cer_beam y wer_beam estan VACIAS para todas las ejecuciones de los
% bloques A-H: la evaluacion con modelo de lenguaje no se relanzo sobre el
% modelo campeon. Es lo unico que falta para cerrar PI3 con datos definitivos.

Los resultados de esta sección proceden de una configuración anterior a la fijación de la receta definitiva y \textbf{deben considerarse preliminares}. Se reportan porque el efecto que documentan es cualitativo y previsiblemente robusto al cambio de codificador, pero sus cifras absolutas habrán de recalcularse.

\begin{table}[H]
  \centering
  \caption{Resultados preliminares de decodificación con modelo de lenguaje de $n$-gramas por fusión superficial.}
  \label{tab:lm-preliminar}
  \small
  \begin{tabular}{lc}
    \toprule
    \textbf{Condición} & \textbf{WER} \\
    \midrule
    Fusión superficial con bigrama & 0,475 \\
    Límite \textit{oracle} de la lista $n$-mejor acústica & 0,503 \\
    \bottomrule
  \end{tabular}
\end{table}

El resultado cualitativamente relevante es que \textbf{la fusión superficial supera al límite \textit{oracle}}. Conforme al argumento de la \cref{sec:oracle}, ningún método de reordenamiento puede superar ese límite, definido como el error que se obtendría eligiendo siempre la mejor hipótesis presente en la lista. Que el sistema quede por debajo demuestra que la transcripción correcta \textbf{no figuraba} en la lista de hipótesis puramente acústicas y que el modelo de lenguaje la alcanzó interviniendo durante la expansión del haz.

La consecuencia para PI3 es directa: una parte del conocimiento lingüístico del sistema completo no reside en el codificador neuronal, y el modelo de lenguaje no es una etapa de posprocesado sino un componente que altera el espacio de búsqueda. El hallazgo se refuerza con el del bloque B: si las unidades de salida gruesas mejoran el texto final en 25 puntos porcentuales, es porque incorporan implícitamente parte de esa misma estructura lingüística dentro del propio codificador. Ambos resultados apuntan en la misma dirección: \textbf{en este sistema, la reconstrucción lingüística pesa más que la lectura de la señal}.

\section{Comparación con el estado del arte}
\label{sec:comparacion-sota}

\begin{table}[H]
  \centering
  \caption{Situación del sistema propuesto frente a la referencia comparable sobre el mismo conjunto de datos.}
  \label{tab:final-vs-sota}
  \small
  \begin{tabular}{llcc}
    \toprule
    \textbf{Sistema} & \textbf{Codificador} & \textbf{CER} & \textbf{WER} \\
    \midrule
    Khanday \textit{et al.} \cite{khanday2026} & Conformer desde cero & 0,238 & --- \\
    Este trabajo, objetivo de carácter & E-Branchformer aleatorio & 0,267 & 0,535 \\
    Este trabajo, objetivo de palabra   & E-Branchformer aleatorio & 0,222 & 0,282 \\
    \bottomrule
  \end{tabular}
\end{table}

Con objetivo de carácter y decodificación voraz, que es la condición directamente comparable, el sistema alcanza $0{,}267$ frente al $0{,}238$ que reportan Khanday \textit{et al.} \cite{khanday2026} sobre el mismo conjunto, la misma unidad y las mismas condiciones de decodificación. La diferencia de unos tres puntos porcentuales \textbf{no es atribuible al preentrenamiento}, puesto que ninguno de los dos sistemas lo emplea en esa configuración, y apunta a la arquitectura del codificador, al aumento de datos o a la receta de optimización.

Es un margen estrecho para un trabajo que no perseguía competir, y sitúa el sistema en el mismo orden de magnitud que la literatura reciente sobre este conjunto. Cambiando el objetivo a palabra la tasa de error por carácter desciende a $0{,}222$, si bien con la salvedad del vocabulario cerrado ya señalada.

La distancia con los sistemas que encabezan la clasificación oficial sigue siendo, no obstante, considerable, y responde a los factores identificados en el \cref{cap:estado-del-arte}: conjuntos de decenas de modelos, modelos de lenguaje de gran tamaño para la fusión de hipótesis y presupuestos de cómputo de otro orden de magnitud.

\section{Discusión}
\label{sec:discusion}

\subsection{Síntesis por pregunta de investigación}

\textbf{PI1, transferencia.} El preentrenamiento de OWSM sobre habla \textbf{no transfiere} a la decodificación de actividad neuronal intracortical. El resultado es reproducible entre semillas, resiste el intento de conceder a la condición preentrenada un ajuste más favorable, y se refuerza con el control decisivo de que un codificador aleatorio congelado iguala o supera al codificador preentrenado congelado con la mitad de presupuesto. Tampoco el tamaño de la arquitectura resulta justificado por la tarea.

\textbf{PI2, granularidad.} La corteza puede leerse a granularidades muy distintas, y la que minimiza el error del texto final es la más gruesa, no la más próxima a la articulación. El factor determinante es la longitud de la secuencia objetivo y la estructura lingüística que la unidad incorpora implícitamente. Este eje produce el mayor efecto medido en todo el trabajo, muy por encima del preentrenamiento.

\textbf{PI3, localización del conocimiento lingüístico.} Una parte sustancial del conocimiento lingüístico reside fuera del codificador. El modelo de lenguaje externo guía la búsqueda en lugar de limitarse a reordenar hipótesis, y las unidades de salida gruesas mejoran el resultado precisamente por internalizar parte de esa misma función.

\subsection{Una lectura conjunta}

Los tres resultados convergen en una conclusión que no formaba parte de las hipótesis iniciales: \textbf{en este problema, y con este volumen de datos, lo que determina el rendimiento no es la calidad de la representación acústica heredada sino la cantidad de estructura lingüística que el sistema logra incorporar}, ya sea a través de la unidad de salida, del modelo de lenguaje externo o de ambos. El preentrenamiento sobre habla aporta lo primero y no lo segundo, lo que ofrece una explicación coherente de por qué su contribución resulta nula.

Frente a ello, el factor limitante del despliegue no es ninguno de los tres ejes estudiados, sino la no estacionariedad: la degradación de 27 puntos porcentuales ante un día de registro no visto supera a cualquier otro efecto medido, y la solución estándar del campo, la capa específica de sesión, resulta ser en ese escenario parte del problema y no de la solución.

\subsection{Amenazas a la validez}

\begin{description}[leftmargin=*]
  \item[Un solo participante.] Todas las conclusiones son específicas de T15, de la localización de sus implantes y del inglés como lengua objetivo.
  \item[Un solo modelo preentrenado.] El resultado se refiere a OWSM. No autoriza a concluir que ningún modelo de habla transfiera, y la evidencia previa sobre wav2vec~2.0 apunta a lo contrario.
  \item[Presupuestos desiguales entre algunas condiciones.] Varias comparaciones enfrentan ejecuciones de 30 y 40 épocas. Cuando esto favorece a la condición que se declara peor, la conclusión resulta conservadora; en el bloque F, en cambio, impide una comparación limpia de eficiencia.
  \item[Condiciones no convergidas.] Las dos condiciones de descongelado parcial conservan pendiente apreciable al final de su presupuesto y sus cifras no son asintóticas.
  \item[Evaluación sobre validación.] Al no distribuirse la partición de prueba con etiquetas, todas las cifras corresponden a validación, sobre la que se han tomado decisiones de diseño, con el consiguiente riesgo de sobreajuste que las cifras no capturan.
  \item[Modelo de lenguaje sin reevaluar.] Las cifras de la \cref{sec:resultados-lm} proceden de una configuración anterior y no se han recalculado sobre el sistema final.
\end{description}
